% BHU PYQ -- Manually transcribed & error-corrected
\documentclass[11pt,a4paper]{article}

\usepackage[a4paper,top=2.5cm,bottom=2.5cm,left=2.8cm,right=2.8cm,headheight=14pt]{geometry}
\usepackage{amsmath,amssymb}
\usepackage[T1]{fontenc}
\usepackage[utf8]{inputenc}
\usepackage{lmodern}
\usepackage{microtype}
\usepackage{fancyhdr}
\usepackage{enumitem}
\usepackage{hyperref}
\usepackage{lastpage}
\usepackage{parskip}

\hypersetup{colorlinks=true,urlcolor=black,linkcolor=black,
  pdftitle={CHB-401: Inorganic & Organic Chemistry-II -- BHU PYQ},pdfauthor={Banaras Hindu University}}

\pagestyle{fancy}\fancyhf{}
\renewcommand{\headrulewidth}{0.4pt}\renewcommand{\footrulewidth}{0.4pt}
\fancyhead[L]{\small\textit{Banaras Hindu University}}
\fancyhead[R]{\small\textit{CHB-401: Inorganic \& Organic Chemistry-II}}
\fancyfoot[C]{\small Page \thepage\ of \pageref{LastPage}}

\newlist{parts}{enumerate}{2}
\setlist[parts,1]{label=(\alph*),leftmargin=*,itemsep=5pt,topsep=3pt}
\setlist[parts,2]{label=(\roman*),leftmargin=*,itemsep=3pt,topsep=2pt}
\newcommand{\pts}[1]{\hfill\textit{[#1]}}

\begin{document}

\begin{center}
  {\large\bfseries BANARAS HINDU UNIVERSITY}\\[2pt]
  {\normalsize B.Sc. (Hons.) Semester IV Examination, 2023-24}\\[6pt]
  \rule{0.8\linewidth}{0.5pt}\\[6pt]
  {\large\bfseries Paper No. CHB-401: Inorganic \& Organic Chemistry-II}\\[4pt]
  {\normalsize Time: 3 Hours\hfill Maximum Marks: 70}
\end{center}
\rule{\linewidth}{0.4pt}
\small
\noindent\textit{Instructions: Answer five questions in total, including Question~1 from each section which is compulsory. Use separate answer books for Section~A and Section~B.}
\normalsize
\bigskip

\begin{center}
  {\large\bfseries SECTION A}\\[2pt]
  {\normalsize (Inorganic Chemistry-II -- Marks: 35)}
\end{center}
\rule{0.4\linewidth}{0.2pt}
\smallskip

\noindent\textbf{Question 1.} Answer all parts of the following: \pts{5 $\times$ 2 = 10}
\begin{parts}
  \item Write all oxidation states shown by the lanthanides $\text{Tb}$ and $\text{Yb}$ with proper reasoning.
  \item Arrange $\text{HBr}$, $\text{HNO}_3$, and $\text{HClO}_4$ according to their acidic strength in: (i) water and (ii) acetic acid. Justify.
  \item Why is anhydrous $\text{CuSO}_4$ colorless and $\text{CuSO}_4\cdot5\text{H}_2\text{O}$ blue? Explain.
  \item Solvents having high dielectric constants are suitable for dissolving ionic solids. Why?
  \item Give two examples of hard and soft acids according to the HSAB principle.
\end{parts}

\medskip\rule{\linewidth}{0.2pt}

\noindent\textbf{Question 2.}
\begin{parts}
  \item Why do $3\text{d}$ transition metals show variable oxidation states? How can their low and high oxidation states be stabilized? Give examples. \pts{6}
  \item State the characteristic features of dilute solutions of sodium in liquid ammonia. How does the electrical conductivity change with increasing concentration of sodium? \pts{6}
\end{parts}

\medskip\rule{\linewidth}{0.2pt}

\noindent\textbf{Question 3.}
\begin{parts}
  \item Write the IUPAC names of the following coordination complexes: \pts{4}
  \begin{parts}
    \item $\text{Na}_3[\text{Co(NO}_2)_6]$
    \item $[\text{Cr(en)}_3]\text{Cl}_3$
    \item $\text{K}[\text{PtCl}_3(\text{C}_2\text{H}_4)]$
  \end{parts}
  \item Draw structures of all possible geometrical and optical isomers of $[\text{Co(en)}_2\text{Cl}_2]^+$. \pts{4}
  \item Write the products of the following reactions in liquid ammonia: \pts{4}
  \begin{parts}
    \item $\text{AgCl} + \text{KNH}_2 \to ?$
    \item $\text{BiI}_3 + 3\text{KNH}_2 \to ?$
  \end{parts}
\end{parts}

\medskip\rule{\linewidth}{0.2pt}

\noindent\textbf{Question 4.}
\begin{parts}
  \item $\text{AsF}_5$ behaves as an acid in $\text{BrF}_3$ solvent. Justify this with chemical equations. \pts{6}
  \item Why is the separation of lanthanides difficult? Describe the ion-exchange method for the separation of lanthanide ions. \pts{6}
\end{parts}

\newpage
\begin{center}
  {\large\bfseries SECTION B}\\[2pt]
  {\normalsize (Organic Chemistry-III -- Marks: 35)}
\end{center}
\rule{0.4\linewidth}{0.2pt}
\smallskip

\noindent\textbf{Question 5.} Answer all parts of the following: \pts{5 $\times$ 2 = 10}
\begin{parts}
  \item Explain why pyridine undergoes electrophilic substitution at position 3, whereas pyridine $1$-oxide reacts at position 4.
  \item Write down the resonating structures of naphthalene and indicate the most stable resonance structure with an explanation.
  \item Fructose has a keto functional group, but it reduces Tollens' and Fehling's reagents. Explain why.
  \item Give the preparation of Indigotin from anthranilic acid.
  \item Compare furan and pyrrole towards electrophilic substitution reactions.
\end{parts}

\medskip\rule{\linewidth}{0.2pt}

\noindent\textbf{Question 6.}
\begin{parts}
  \item Write down the structure of the product of each of the following reactions: \pts{2 $\times$ 3 = 6}
  \begin{parts}
    \item $\text{Naphthalene} \xrightarrow{\text{alkaline KMnO}_4} ?$
    \item $\text{Indigotin} \xrightarrow{\text{dil. NaOH, Zn}} ?$
    \item $\text{Pyrrole} \xrightarrow{\text{CH}_3\text{MgBr, then CH}_3\text{I}} ?$
  \end{parts}
  \item Complete the following reactions, showing the mechanisms: \pts{6}
  \begin{parts}
    \item $\text{Pyridine} \xrightarrow{\text{NaNH}_2, \text{toluene}} ?$ (Chichibabin reaction)
    \item $\text{Furan} \xrightarrow{\text{CH}_3\text{COCl, SnCl}_4} ?$
  \end{parts}
\end{parts}

\medskip\rule{\linewidth}{0.2pt}

\noindent\textbf{Question 7.}
\begin{parts}
  \item Write down the structures of the products of each of the following reactions: \pts{2 $\times$ 3 = 6}
  \begin{parts}
    \item $\beta\text{-D-Glucose} \xrightarrow{\text{dry HCl, MeOH}} ?$
    \item $\text{Quinoline} \xrightarrow{\text{HNO}_3, \text{H}_2\text{SO}_4} ?$
    \item $\text{Anthracene} \xrightarrow{\text{Na, EtOH}} ?$
  \end{parts}
  \item Discuss Skraup quinoline synthesis with a detailed mechanism. \pts{6}
\end{parts}

\medskip\rule{\linewidth}{0.2pt}

\noindent\textbf{Question 8.}
\begin{parts}
  \item How will you carry out the following conversions? Write stepwise mechanisms: \pts{6}
  \begin{parts}
    \item $\text{D-Glucose} \xrightarrow{\text{i. Br}_2/\text{H}_2\text{O, ii. Fe}^{3+}, \text{H}_2\text{O}_2} ?$ (Ruff degradation)
    \item $\text{D-Arabinose} \xrightarrow{\text{i. HCN, ii. H}_3\text{O}^+, \text{iii. Na-Hg}} ?$ (Kiliani-Fischer synthesis)
  \end{parts}
  \item Discuss the structures and aromatic characters of pyrrole, furan, and thiophene. \pts{6}
\end{parts}

\vfill
\rule{\linewidth}{0.4pt}
\begin{center}\small
  \href{https://bhu-exam.vercel.app/}{Click here to visit our website}\\[4pt]
  \href{https://me-aryan.vercel.app/}{Click here to visit our contributor}\\[4pt]
  \href{https://quashan-me.vercel.app/}{Click here to visit our contributor}
\end{center}

\end{document}
